\chapter*{Abstract}

% Time series? *!
Time series data are ubiquitous, appearing inherently in diverse fields or as transformations of non-temporal data such as biological sequences (strings).
%
Many data mining tasks are defined on them, such as searching (query-by-content), motif (nearest-neighbor pair) finding, transformation (feature extraction), and classification.
%
Given the prosperity of time series analysis, the development of robust and more flexible analytical frameworks is essential.

% One-sentence description *!
We first provide a one-sentence description of this thesis.
%
This thesis develops three model-agnostic, subsequence-centric methods to advance time series analysis—by 
proposing efficient piecewise scaling for more flexible similarity computation, 
generating subsequence-derived covariates to improve forecasting accuracy, 
and encoding domain-specific biological subsequences into multivariate time series for classification—demonstrating 
how targeted exploitation of local patterns can benefit diverse downstream tasks from general time series modeling to bioinformatics.
%
We then explain the three methods part by part.

% Study 1 *!
The first part addresses the limitations of current similarity measures.
%
% Similarity search is a core subroutine in time series data mining tasks.
While Dynamic Time Warping (DTW) and Uniform Scaling (US) are prevailing measures for handling local distortions and global scaling, respectively, and some studies have demonstrated that combining both DTW and US is necessary to obtain meaningful results.
However, the current approaches fail when different phases of a process are expressed at varying speeds. They assume a single scaling factor for the entire sequence. 
%
We introduce Piecewise Scaling Distance (PSD), the first framework that achieves invariance to multiple local scaling factors, and incorporates efficient speed-up techniques. 
PSD provides a more flexible, expressive, and interpretable measure of similarity, as demonstrated by its ability to segment and identify different scaling factors within the two compared time series during the computation of the similarity measure.

% Study 2 *!
The second part exploits a data mining primitive, namely Matrix Profile, to enhance time series forecasting. 
In brief, Matrix Profile encodes the information about the nearest neighbor for each subsequence of a time series.
We argue that historical subsequence occurrences contain predictive power. 
By identifying the $k$-nearest neighbors of subsequences in the past using Matrix Profile, we extract historical occurrence information as covariates for a forecaster. 
Our results demonstrate that these computationally efficient, data-driven covariates significantly improve forecasting accuracy across various domains.

% Study 3 *!
The last part demonstrates the utility of time series analysis in bioinformatics through the prediction of Human Dicer Cleavage Sites. 
%
Current methods for RNA sequence analysis often rely on complex feature engineering or computationally expensive deep learning models. 
We propose MTSCCleav, a method that encodes 1-D RNA sequences and base-pair probabilities in predicted secondary structures into multivariate time series.
To the best of our knowledge, we are the first to make use of the base-pair probabilities in predicted secondary structures.
This approach allows the application of well-established time series tools to biological problems. 
Experiments show that MTSCCleav achieves state-of-the-art accuracy with a 3.7x to 28.8x speedup, while perturbation analysis confirms that pre-miRNA center regions are critical for cleavage-site prediction, consistent with the existing literature.

% Conclusion
Collectively, this thesis advances the field of time series analysis by demonstrating that targeted, subsequence-level insights—through piecewise scaling, covariate augmentation, and encoding—provide a unified and powerful perspective for addressing fundamental challenges from general data miningto bioinformatics.

% \noindent\rule{\linewidth}{0.4pt}
%%
%%
%%






